% Full title: M\"uller Representation and the Coupled Operator
\section{M\"uller Representation and the Coupled Operator}
\label{sec:muller-coupling}

\subsection{Rayleigh and layer-potential conventions}

Let $\nu$ denote the unit normal on $\Upsilon^0=\partial\Omega^0$ pointing from
$\Omega^0$ into the background.  Interior and exterior traces are
\[
  \gamma^{\mathrm i}u(z):=\lim_{h\downarrow0}u(z-h\nu(z)),
  \qquad
  \gamma^{\mathrm e}u(z):=\lim_{h\downarrow0}u(z+h\nu(z)).
\]
The same superscripts are used for normal derivatives in the direction $\nu$.
The outgoing two-dimensional fundamental solution is
\begin{equation}
  \Phi_\kappa(x,z):=\frac{\ii}{4}H_0^{(1)}(\kappa|x-z|).
  \label{eq:fundamental-solution}
\end{equation}
For $x\notin\Upsilon^0$, define
\begin{align}
  (S_\kappa\sigma)(x)&:=\int_{\Upsilon^0}\Phi_\kappa(x,z)\sigma(z)\dd s(z),
  \label{eq:single-layer}\\
  (D_\kappa\tau)(x)&:=\int_{\Upsilon^0}
  \partial_{\nu(z)}\Phi_\kappa(x,z)\tau(z)\dd s(z).
  \label{eq:double-layer}
\end{align}
With these conventions,
$\gamma_D^{\mathrm e}D_\kappa-\gamma_D^{\mathrm i}D_\kappa=I$ and
$\gamma_N^{\mathrm e}S_\kappa-\gamma_N^{\mathrm i}S_\kappa=-I$.

Let $G^{\mathrm{qp}}_{k,\beta}$ be the vertically $\beta$-quasiperiodic
outgoing Green function satisfying
\[
  G^{\mathrm{qp}}_{k,\beta}(x+de_y,z)
  =e^{\ii\beta d}G^{\mathrm{qp}}_{k,\beta}(x,z),
  \qquad (\Delta_x+k^2)G^{\mathrm{qp}}_{k,\beta}(x,z)=-\delta_z.
\]
Away from its exceptional set it is represented by the lattice sum
\begin{equation}
  G^{\mathrm{qp}}_{k,\beta}(x,z)
  =\sum_{\ell\in\Z}e^{\ii\beta\ell d}
  \Phi_k(x,z+\ell d e_y),
  \label{eq:qp-Green}
\end{equation}
interpreted by the standard convergent continuation when necessary.  Replacing
$\Phi_k$ by $G^{\mathrm{qp}}_{k,\beta}$ in
\eqref{eq:single-layer}--\eqref{eq:double-layer} defines
$S^{\mathrm{qp}}_{k,\beta}$ and $D^{\mathrm{qp}}_{k,\beta}$.

For
\[
  \gamma_m(k,\beta):=\sqrt{k^2-\beta_m^2},
  \qquad \operatorname{Im}\gamma_m\geq0,
\]
a Wood anomaly occurs when $\gamma_m=0$ for some $m$.  At a non-Wood
parameter, define
\begin{equation}
  h^{1/2}_\beta:=
  \left\{a=(a_m)_{m\in\Z}:
  \sum_m\langle\beta_m\rangle|a_m|^2<\infty\right\},
  \qquad
  \Xi_\beta:=h^{1/2}_\beta\times h^{1/2}_\beta.
  \label{eq:Rayleigh-space}
\end{equation}
For $\xi=(\xi^\rightarrow,\xi^\leftarrow)\in\Xi_\beta$, the homogeneous
Rayleigh reconstruction in the center cell is
\begin{equation}
  (\mathcal H_{k,\beta}\xi)(x,y)
  :=\sum_{m\in\Z}\left[
  \xi_m^\rightarrow e^{\ii\gamma_m(x-X^-)}
  +\xi_m^\leftarrow e^{-\ii\gamma_m(x-X^+)}\right]\psi_m(y).
  \label{eq:Rayleigh-reconstruction}
\end{equation}
The weights in \eqref{eq:Rayleigh-space} control both the
$H^{1/2}_\beta$ Dirichlet trace and the $H^{-1/2}_\beta$ normal trace.

The interface-density space is
\begin{equation}
  \mathcal D^0:=H^{1/2}(\Upsilon^0)\times H^{-1/2}(\Upsilon^0),
  \qquad \eta=(\tau,\sigma)\in\mathcal D^0,
  \label{eq:density-space}
\end{equation}
where $\tau$ is the double-layer density and $\sigma$ is the single-layer
density.  To describe a field before its two material traces have been
matched, set
\begin{equation}
  H^1_{\beta,\mathrm{br}}(\mathcal C^0):=
  H^1_\beta(\mathcal C^0\setminus\overline{\Omega^0})
  \oplus H^1(\Omega^0).
  \label{eq:center-broken-space}
\end{equation}
The subscript ``br'' means \emph{broken across $\Upsilon^0$}: the two components
have independent traces until the transmission residual below is set to zero.
Let $k^0:=n^0k$.  The center reconstruction
\begin{equation}
  \mathcal R_{\mathrm c}(k,\beta):\mathcal D^0\times\Xi_\beta
  \longrightarrow H^1_{\beta,\mathrm{br}}(\mathcal C^0)
  \label{eq:center-reconstruction-map}
\end{equation}
is the broken field $u_{\eta,\xi}^{\mathrm c}$ with
\begin{equation}
  u_{\eta,\xi}^{\mathrm c}(x)=
  \begin{cases}
  \mathcal H_{k,\beta}\xi(x)
  +D^{\mathrm{qp}}_{k,\beta}\tau(x)
  -S^{\mathrm{qp}}_{k,\beta}\sigma(x),
  &x\in\mathcal C^0\setminus\overline{\Omega^0},\\
  D_{k^0}\tau(x)-S_{k^0}\sigma(x),
  &x\in\Omega^0.
  \end{cases}
  \label{eq:center-reconstruction}
\end{equation}
Thus the exterior component is vertically quasiperiodic, while the bounded
interior component needs no separate quasiperiodic condition.

Define the M\"uller transmission residual by taking all traces in the fixed
normal direction $\nu$:
\begin{equation}
\begin{split}
  \mathcal M(k,\beta)(\eta,\xi):=
  \bigl(&\gamma_D^{\mathrm e}u_{\eta,\xi}^{\mathrm c}
         -\gamma_D^{\mathrm i}u_{\eta,\xi}^{\mathrm c},\\
        &\gamma_N^{\mathrm e}u_{\eta,\xi}^{\mathrm c}
         -\gamma_N^{\mathrm i}u_{\eta,\xi}^{\mathrm c}\bigr)
  \in\mathcal H_{\Upsilon^0},
  \qquad
  \mathcal H_{\Upsilon^0}:=H^{1/2}(\Upsilon^0)\times H^{-1/2}(\Upsilon^0).
\end{split}
\label{eq:Muller-residual}
\end{equation}
This trace definition fixes every half-identity sign without relying on an
unwritten matrix convention.  Its expansion using the jump relations is the
M\"uller block.  Let
\begin{equation}
  \mathcal Z_{\mathrm c}(k,\beta):=\ker\mathcal M(k,\beta),
  \qquad
  \mathcal N_{\mathrm c}(k,\beta):=
  \mathcal Z_{\mathrm c}(k,\beta)\cap\ker\mathcal R_{\mathrm c}(k,\beta).
  \label{eq:center-nullspaces}
\end{equation}
Thus $\mathcal N_{\mathrm c}$ contains exactly the center algebraic
representations that satisfy the transmission equations but generate the zero
physical center field.

Let $\mathcal S_{\mathrm c}(k,\beta)$ be the space of piecewise Helmholtz
fields in $H^1_\beta(\mathcal C^0)$ satisfying the physical TM transmission
conditions at $\Upsilon^0$, with no boundary condition at $\Gamma^\pm$.
Define the representation exceptional set $\mathfrak E_{\mathrm{rep}}(\beta)$
as the union of
\begin{enumerate}[label=(\roman*),leftmargin=2.2em]
  \item Wood parameters $\gamma_m(k,\beta)=0$;
  \item poles or undefined parameters of the chosen quasiperiodic Green
  function;
  \item parameters for which the complementary swapped-wavenumber transmission
  problem with interior wavenumber $k$ and exterior wavenumber $k^0$ is not
  uniquely solvable.
\end{enumerate}
The quotient below retains possible representation nullspaces even outside
this explicitly listed set.

\begin{theorem}[Center-cell M\"uller--Rayleigh representation]
\label{thm:center-representation}
Let $k^2\in I_\beta$ and $k\notin\mathfrak E_{\mathrm{rep}}(\beta)$.  The
restriction of $\mathcal R_{\mathrm c}(k,\beta)$ to
$\mathcal Z_{\mathrm c}(k,\beta)$ is surjective onto
$\mathcal S_{\mathrm c}(k,\beta)$.  Consequently it induces a linear
isomorphism
\begin{equation}
  \mathcal Z_{\mathrm c}(k,\beta)/\mathcal N_{\mathrm c}(k,\beta)
  \xrightarrow{\ \cong\ }\mathcal S_{\mathrm c}(k,\beta),
  \qquad [\eta,\xi]\longmapsto\mathcal R_{\mathrm c}(k,\beta)(\eta,\xi).
  \label{eq:center-quotient-isomorphism}
\end{equation}
No uniqueness of the densities is asserted.  Density uniqueness is equivalent
to the additional statement $\mathcal N_{\mathrm c}(k,\beta)=\{0\}$.
\end{theorem}
\status{adaptation and proof target.}
\begin{proofroadmap}
Extract from a given center field the two homogeneous Rayleigh components in
homogeneous collars of $\Gamma^-$ and $\Gamma^+$.  Subtract their continuation
\eqref{eq:Rayleigh-reconstruction}; the remainder satisfies the outgoing
single-cell representation hypotheses.  Apply the quasiperiodic exterior and
free-space interior Green identities and construct the complementary
swapped-wavenumber field as in
\cite[Theorem~4 and Appendix~A, Lemmas~9--10]{BarnettGreengard2010}.  The
exterior double layer in this calculation is
$D^{\mathrm{qp}}_{k,\beta}$, whereas the inclusion double layer is $D_{k^0}$;
using $D^{\mathrm{qp}}_{k^0,\beta}$ in the exterior extinction identity is the
kernel mismatch that must not be repeated.

Surjectivity of the physical-field map is the substantive claim.  Once it is
proved, the quotient isomorphism follows from the first isomorphism theorem.
The density kernel must be analyzed separately with jump relations,
complementary fields, and Calder\'on projector ranges; the transmission-BIE
analysis in \cite[Lemmas~2.2--2.6]{HiptmairMoiolaSpence2022} explains why a
physical uniqueness theorem does not automatically prove density uniqueness.

Detailed internal guidance is in
\path{research/planning/muller-cauchy/notes/layers.md} and the
obligations named ``corrected center representation'' and
``density-to-field kernel characterization'' in
\path{research/planning/muller-cauchy/p-proofs.md}.  The weakest
acceptable result is surjectivity modulo a finite-dimensional missed
homogeneous sector; an unquotiented theorem is acceptable only after
$\mathcal N_{\mathrm c}=\{0\}$ is proved on the stated regular set.
\end{proofroadmap}
\begin{proof}
Write $u=(u^{\mathrm e},u^{\mathrm i})\in
\mathcal S_{\mathrm c}(k,\beta)$ for the exterior and inclusion components,
and construct $(\eta,\xi)\in\mathcal Z_{\mathrm c}(k,\beta)$ whose
reconstructed field is $u$.

\emph{Step 1: separate the incoming homogeneous Rayleigh field.}
Because $\Omega^0\Subset\mathcal C^0$, there are homogeneous collars of both
ports.  In either collar, the Fourier coefficient of $u^{\mathrm e}$ in the
mode $\psi_m$ solves
\[
  \partial_x^2u_m+\gamma_m^2u_m=0.
\]
The non-Wood assumption gives $\gamma_m\neq0$, so the solution has a unique
right/left-going decomposition.  Denote by $a_m^{-,\rightarrow}$ the
right-going coefficient in the left collar and by
$a_m^{+,\leftarrow}$ the left-going coefficient in the right collar, with
the phases based respectively at $X^-$ and $X^+$.  If $(f_m^\pm,p_m^\pm)$
are the Fourier coefficients of the Dirichlet and $\partial_x$ traces at the
two ports, then
\[
  a_m^{-,\rightarrow}
  =\frac12\left(f_m^-+\frac{p_m^-}{\ii\gamma_m}\right),
  \qquad
  a_m^{+,\leftarrow}
  =\frac12\left(f_m^+-\frac{p_m^+}{\ii\gamma_m}\right).
\]
The trace theorem gives $f^\pm\in h^{1/2}_\beta$ and
$p^\pm\in h^{-1/2}_\beta$.  Moreover,
$|\gamma_m|\asymp\langle\beta_m\rangle$ for all sufficiently large $|m|$;
the remaining set is finite and contains no zero $\gamma_m$.  Hence
\[
  \xi:=\bigl((a_m^{-,\rightarrow})_m,
             (a_m^{+,\leftarrow})_m\bigr)\in\Xi_\beta.
\]
Set $h:=\mathcal H_{k,\beta}\xi$ and $w:=u^{\mathrm e}-h$.  In the left
collar $w$ contains only left-going modes, and in the right collar it contains
only right-going modes: the term launched from the opposite port contributes
only to the corresponding outgoing component.  Continuing these modal
series through the two ports therefore extends $w$ to a vertically
$\beta$-quasiperiodic solution of $(\Delta+k^2)w=0$ in
$B\setminus\overline{\Omega^0}$ which is outgoing at both ends.

\emph{Step 2: the two Green representation identities.}
Use the fixed normal $\nu$ and abbreviate
\[
  P_\kappa(f,q):=D_\kappa f-S_\kappa q,
  \qquad
  P^{\mathrm{qp}}_{k,\beta}(f,q)
  :=D^{\mathrm{qp}}_{k,\beta}f-S^{\mathrm{qp}}_{k,\beta}q.
\]
Green's second identity, first on a truncated quasiperiodic strip and then by
the outgoing limiting-absorption limit, gives
\[
  P^{\mathrm{qp}}_{k,\beta}
  (\gamma_D^{\mathrm e}w,\gamma_N^{\mathrm e}w)
  =w
  \quad\text{in }B\setminus\overline{\Omega^0}.
\]
The same identity applied to a solution $z$ of $(\Delta+k^2)z=0$ in
$\Omega^0$ yields the extinction relation
\[
  P^{\mathrm{qp}}_{k,\beta}
  (\gamma_D^{\mathrm i}z,\gamma_N^{\mathrm i}z)=0
  \quad\text{in }B\setminus\overline{\Omega^0}.
\]
For the free-space kernel at $k^0$, the corresponding identities are
\begin{align*}
  P_{k^0}(\gamma_D^{\mathrm i}z,\gamma_N^{\mathrm i}z)
  &=-z&&\text{in }\Omega^0,\\
  P_{k^0}(\gamma_D^{\mathrm e}z,\gamma_N^{\mathrm e}z)
  &=0&&\text{in }\Omega^0
\end{align*}
when the latter $z$ is an outgoing exterior field.  The signs follow from
the jump conventions preceding \eqref{eq:single-layer}.  These are the
one-periodic analogues of the identities used in
\cite[Appendix~A, Lemmas~9--10]{BarnettGreengard2010}; here the vertical
wall terms cancel by $\beta$-quasiperiodicity and the two end terms are
removed by outgoing limiting absorption.

\emph{Step 3: construct common M\"uller densities.}
Let
\[
  f:=\gamma_D^{\mathrm e}u^{\mathrm e}
    =\gamma_D^{\mathrm i}u^{\mathrm i},
  \qquad
  q:=\gamma_N^{\mathrm e}u^{\mathrm e}
    =\gamma_N^{\mathrm i}u^{\mathrm i},
\]
and write $h_D:=\gamma_Dh$, $h_N:=\gamma_Nh$ on $\Upsilon^0$.  Then the
exterior Cauchy data of $w$ are $(f-h_D,q-h_N)$.  Consider the complementary
swapped-wavenumber transmission problem
\begin{align*}
  (\Delta+k^2)v^{\mathrm i}&=0
    &&\text{in }\Omega^0,\\
  (\Delta+(k^0)^2)v^{\mathrm e}&=0
    &&\text{in }\R^2\setminus\overline{\Omega^0},
\end{align*}
with the Sommerfeld condition for $v^{\mathrm e}$ and jumps
\begin{align*}
  \gamma_D^{\mathrm e}v^{\mathrm e}
  -\gamma_D^{\mathrm i}v^{\mathrm i}&=2f-h_D,\\
  \gamma_N^{\mathrm e}v^{\mathrm e}
  -\gamma_N^{\mathrm i}v^{\mathrm i}&=2q-h_N.
\end{align*}
Clause (iii) in the definition of
$\mathfrak E_{\mathrm{rep}}(\beta)$ and
$k\notin\mathfrak E_{\mathrm{rep}}(\beta)$ give its unique solution for
these admissible trace data.  Define
\begin{align*}
  \tau
  &:=f-h_D+\gamma_D^{\mathrm i}v^{\mathrm i}
    =-f+\gamma_D^{\mathrm e}v^{\mathrm e},\\
  \sigma
  &:=q-h_N+\gamma_N^{\mathrm i}v^{\mathrm i}
    =-q+\gamma_N^{\mathrm e}v^{\mathrm e}.
\end{align*}
The equalities are exactly the two prescribed jumps, and the trace mapping
properties put $(\tau,\sigma)$ in $\mathcal D^0$.

Set $\eta=(\tau,\sigma)$.  By the two quasiperiodic identities above, the
exterior component of the reconstruction is
\begin{align*}
  h+P^{\mathrm{qp}}_{k,\beta}(\tau,\sigma)
  &=h+P^{\mathrm{qp}}_{k,\beta}(f-h_D,q-h_N)\\
  &\quad
    +P^{\mathrm{qp}}_{k,\beta}
      (\gamma_D^{\mathrm i}v^{\mathrm i},
       \gamma_N^{\mathrm i}v^{\mathrm i})\\
  &=h+w=u^{\mathrm e}.
\end{align*}
Likewise, using the alternative expressions for $\tau,\sigma$ and the two
free-space identities above, its interior component is
\begin{align*}
  P_{k^0}(\tau,\sigma)
  &=-P_{k^0}(f,q)
    +P_{k^0}(\gamma_D^{\mathrm e}v^{\mathrm e},
             \gamma_N^{\mathrm e}v^{\mathrm e})\\
  &=u^{\mathrm i}.
\end{align*}
Thus $\mathcal R_{\mathrm c}(k,\beta)(\eta,\xi)=u$.  Since $u$ satisfies
the physical transmission conditions, its two jumps vanish, so
$\mathcal M(k,\beta)(\eta,\xi)=0$ and
$(\eta,\xi)\in\mathcal Z_{\mathrm c}(k,\beta)$.  This proves the asserted
surjectivity.

\emph{Step 4: pass to the quotient.}
The kernel of
$\mathcal R_{\mathrm c}|_{\mathcal Z_{\mathrm c}}$ is, by
\eqref{eq:center-nullspaces}, exactly $\mathcal N_{\mathrm c}$.  The first
isomorphism theorem therefore gives
\eqref{eq:center-quotient-isomorphism}.  This argument neither requires nor
asserts that $\mathcal N_{\mathrm c}$ is zero.
\end{proof}

\subsection{The continuous coupled operator}

The port trace of the center reconstruction is the bounded map
\begin{equation}
  \Pi^\pm(k,\beta):=\Tr^{0,\pm}\mathcal R_{\mathrm c}(k,\beta):
  \mathcal D^0\times\Xi_\beta\longrightarrow\mathcal H_\Gamma^\pm.
  \label{eq:center-port-map}
\end{equation}
This definition includes both the layer contribution and the homogeneous
Rayleigh contribution.  The half-guide maps $\mathcal E^\pm$ and $Q^\pm$ were
defined in \eqref{eq:field-synthesis}--\eqref{eq:trace-synthesis}.

\begin{definition}[Coupled M\"uller--Cauchy operator and representation nullspace]
\label{def:coupled-operator}
Define the algebraic domain and residual codomain by
\begin{align}
  \mathcal X(k,\beta)
  &:=\mathcal D^0\times\Xi_\beta\times
  \mathcal K_s^-\times\mathcal K_s^+,
  \label{eq:coupled-domain}\\
  \mathcal Y(k,\beta)
  &:=\mathcal H_{\Upsilon^0}\times
  \mathcal H_\Gamma^-\times\mathcal H_\Gamma^+.
  \label{eq:coupled-codomain}
\end{align}
For $x=(\eta,\xi,c^-,c^+)\in\mathcal X(k,\beta)$, set
\begin{equation}
  \mathbb A(k,\beta)x:=
  \begin{bmatrix}
  \mathcal M(k,\beta)(\eta,\xi)\\
  \Pi^-(k,\beta)(\eta,\xi)-Q^-(k,\beta)c^-\\
  \Pi^+(k,\beta)(\eta,\xi)-Q^+(k,\beta)c^+
  \end{bmatrix}
  \in\mathcal Y(k,\beta).
  \label{eq:coupled-operator}
\end{equation}
The first row is the physical transmission mismatch on $\Upsilon^0$; the second
and third rows are the complete center-oriented Cauchy mismatches at the left
and right ports.

Let
\begin{equation}
  H^1_{\beta,\mathrm{br}}(B):=
  H^1_\beta(W^- )\oplus
  H^1_\beta(\mathcal C^0\setminus\overline{\Omega^0})
  \oplus H^1(\Omega^0)
  \oplus H^1_\beta(W^+).
  \label{eq:global-broken-space}
\end{equation}
The global broken-field reconstruction is
\begin{equation}
  \mathcal R_{\mathrm g}(k,\beta)x:=
  \begin{cases}
    \mathcal E^-(k,\beta)c^-,&\text{in }W^-,\\
    \mathcal R_{\mathrm c}(k,\beta)(\eta,\xi),&\text{in }\mathcal C^0,\\
    \mathcal E^+(k,\beta)c^+,&\text{in }W^+.
  \end{cases}
  \label{eq:global-reconstruction}
\end{equation}
It is a broken field for general $x$ and a global $H^1_\beta(B)$ field when
$x\in\ker\mathbb A(k,\beta)$.  Finally, the representation-induced nullspace
is
\begin{equation}
  \mathcal N_{\mathrm{rep}}(k,\beta):=
  \left\{x\in\ker\mathbb A(k,\beta):
  \mathcal R_{\mathrm g}(k,\beta)x=0
  \text{ in }H^1_{\beta,\mathrm{br}}(B)\right\}.
  \label{eq:Nrep-definition}
\end{equation}
Thus $\mathbb A$ is the fully defined algebraic coupling, while
$\mathcal N_{\mathrm{rep}}$ records precisely its kernel vectors that carry no
physical field.  Neither symbol is shorthand for an unspecified numerical
matrix.
\end{definition}
